\section{A CM local boundary and its global projective-image obstruction}
\label{cb-section}

The boundary considered here differs from the non-CM seed $(0,1)$.
It supplies actual positive seeds satisfying a specified finite collection
of local tests, while a separate global theorem excludes congruence to
CM newforms. Both assertions concern their exact stated tests and
representations.

Write $N(a,b)=a^2+ab+b^2$ and retain the actual quartic $F(a,b)$
and Frey curve $E_{a,b}$ defined in the preceding modular sections.

\begin{lemma}[An explicit CM boundary]\label{cb-cm-boundary}
At $(a,b)=(-1,1)$ one has $N(a,b)=F(a,b)=1$, and the actual Frey
curve specializes over $K=\mathbb Q(r)$, $r^2=-3$, to
\[
 E_*:\quad Y^2=X^3+12rX.
\]
It has geometric CM by $\mathbb Z[i]$, good reduction away from the
primes over $6$, and trace zero at both places over every rational
prime $q\equiv7\pmod{12}$.
\end{lemma}
\begin{proof}
The specialization follows by substitution. Its discriminant is
$-64(12r)^3$, and its $j$-invariant is $1728$. Over $K(i)$ the
origin-preserving automorphism $(X,Y)\mapsto(-X,iY)$ squares to
elliptic negation, giving the stated endomorphisms.
For $q\equiv7\pmod{12}$, the prime splits in $K$ and $-1$ is a
quadratic nonsquare modulo $q$. At either place choose $r_q^2=-3$.
The polynomial $f(X)=X^3+12r_qX$ is odd, so the quadratic-character
sum of $f(X)$ cancels in pairs $X,-X$, with zero contribution at
$X=0$. Thus the full point count is $q+1$ and its trace is zero.
\end{proof}

\begin{theorem}[Actual positive finite local shadows]\label{cb-local-shadows}
For every positive integer $L$ and every height target, there exist
positive coprime $a,b$ of greater height satisfying
\[
 (a,b)\equiv(-1,1)\pmod L,
 \qquad N(a,b)\equiv F(a,b)\equiv1\pmod L.
\]
Their Frey coefficients agree with those of $E_*$ modulo $L$.
Consequently they satisfy all direct norm power-residue tests of any
positive exponent at divisors of $L$, and all the trace-zero tests
at primes $q\mid L$ with $q\equiv7\pmod{12}$.
\end{theorem}
\begin{proof}
Take $b=1$ and $a=Lk-1$ for arbitrarily large positive integers $k$.
Primitivity is immediate and $a+b=Lk$ tends to infinity. Evaluate the
two integer polynomials at $(-1,1)$ to get the norm congruences;
the root $1$ certifies every stated power-residue test. The polynomial
Frey coefficients also specialize modulo $L$. Their reductions at
the stated primes are good and identical to $E_*$, so their traces
are zero by Lemma~\ref{cb-cm-boundary}.
Taking $L$ divisible by specified prime powers preserves any specified
finite coefficient precision. The mixed seed remains primitive because
a prime dividing both $ab$ and $N$ would divide both $a,b$; its square
identities hold exactly for these actual integer seeds.
\end{proof}

These shadows assert no global power equation. In fact
Theorem~\ref{sgap-axis} excludes a square $F$ throughout this family,
since $b=1$. Theorem~\ref{as-finite} makes all global pure powers
$F(a,1)=Q^g$, $Q>1$, $g\ge2$, effectively finite across all
exponents. It does not classify the remaining odd-exponent possibilities.
Nor does the local shadow identify a particular CM newform or supply
an isomorphism of residual Galois modules.

\begin{theorem}[Global exclusion of CM residual candidates]\label{cb-cm-exclusion}
Let $a,b$ be positive coprime integers and $p>7$ a prime. The specified
residual descent of $E_{a,b}[p]\otimes\chi$ used in the Frey modular
reduction is not isomorphic, over an algebraic closure of its residue
coefficient field, to the semisimplified residual representation of
any characteristic-zero CM newform.
\end{theorem}
\begin{proof}
The actual quartic satisfies $F\ge13$ and $\gcd(F,6)=1$, so some
prime $q>3$ divides it. The proved local calculation makes $E_{a,b}$
multiplicative at the places over $q$, and its explicit cyclic
isogeny to its conjugate has square-free degree two. The general
Q-curve large-image theorem, with the same published $N_3=7$ input
already checked in the Frey modular section, therefore applies.

For precision about its target, Ellenberg's
\href{https://people.math.wisc.edu/~ellenberg/A4B2Cp.pdf}{author-hosted original paper},
Theorem~3.14, displays a surjective projective extension
\[
 \mathbb P\bar\rho_{E,p}:G_{\mathbb Q}\longrightarrow
       \operatorname{PGL}_2(\mathbb F_p),
\]
unless the curve has potentially good reduction at all primes outside
$6$. The multiplicative place excludes that alternative. The explicit
$N_3=7$ specialization is the general Theorem~5.2 in
\href{https://sweet.ua.pt/apacetti/papers/Q-curves.pdf}{Pacetti--Villagra Torcomian};
it is an external published input, not a new computation or an inference
from a nonprimitive substitution in another Diophantine equation.

On $G_K$, the image $H$ of this projective representation is a normal
subgroup of index at most two in $\operatorname{PGL}_2(\mathbb F_p)$.
It contains $\operatorname{PSL}_2(\mathbb F_p)$. Indeed the latter
is perfect for $p>3$: $\operatorname{SL}_2(\mathbb F_p)$ is generated
by upper and lower elementary unipotents, and, for $t^2\ne1$,
\[
 [\operatorname{diag}(t,t^{-1}), U(x/(t^2-1))]=U(x).
\]
The analogous lower formula proves perfectness before and after
projectivization. Every abelian quotient of the projective group kills
that subgroup. Its projective image is nontrivial, so it is not
solvable; neither is $H$. We have not asserted that the image on
$G_K$ itself must be the full projective linear group.

A CM newform's characteristic-zero representation is induced from a
one-dimensional character of the Galois group of its imaginary
quadratic CM field. This CM representation description is also used
explicitly in Section~6 of the same primary Q-curve paper. In an
induced stable lattice its two cosets give diagonal matrices on the
index-two subgroup and antidiagonal matrices on the other coset.
After reduction, if the two inducing characters differ, the resulting
module is irreducible of this monomial form; if they coincide, it
splits over $\overline{\mathbb F}_p$, since $p$ is odd and the
coset matrix has two distinct square-root eigenvalues. Thus the
semisimplified projective image is abelian or cyclic-by-a-group-of-order
at-most-two, and is solvable. The diagonal projective image is finite
cyclic as a finite subgroup of $\overline{\mathbb F}_p^*$.

Any proposed residual isomorphism can now be restricted to $G_K$.
Projectivization removes the scalar character $\chi$. The actual side
has the nonsolvable image $H$, whereas the CM side has a subgroup
of a solvable group. This is a contradiction. Coefficient embeddings
and scalar twists do not affect it. Large image also ensures that
semisimplifying the actual module loses no extension in this comparison.
\end{proof}

This uses uniform large projective image. The fact that an elliptic
curve individually has no geometric CM would not suffice to exclude
all its CM residual congruences at the stated primes.

\begin{corollary}[Two remaining pure-branch orbits]\label{cb-two-orbits}
If an actual seed satisfies $F=Q^p$ with $Q>1$ and prime $p>7$,
its fixed-level modular candidate belongs to one of the two non-CM
Galois orbits at levels $288$ and $576$ in the complete exact
newspace enumeration. All coefficient embeddings remain candidates.
\end{corollary}
\begin{proof}
The pure-branch modular theorem supplies weight two, character
$\psi_3$ and one of the five levels $36,72,144,288,576$.
The independently repeated complete enumeration in the preceding
manuscript lists precisely two non-CM orbits. Apply
Theorem~\ref{cb-cm-exclusion} to all the CM candidates.
\end{proof}

The same projective-image argument removes CM candidates at the
varying levels and non-flat weights of the other modular reductions.
Those spaces are not reduced to the same two-orbit list. The two
non-CM pure-branch orbits are not excluded here, and the level-$576$
boundary comparison continues to retain the distinction between
$\chi$ and $\chi^{-1}$.

\paragraph{Verification and remaining scope.}
All five ordinary claims passed complete independent reviews by the
other two research agents; both reopened the original projective-image
source and the published small-prime input. A separate exact replay
enumerates all $81{,}376$ affine states in the sixteen listed finite
reductions and checks thirty-two actual positive congruence shadows.
These finite checks do not establish the global representation theorem.
No new Lean proof of projective-image surjectivity, CM induction, or
modular elimination is asserted. The finite local shadows and the
global CM obstruction coexist: the latter closes that specific
candidate route while leaving the non-CM modular and analytic tail
routes open.
